\section{Discussion}\label{sec:discussion}

\subsection{Why Magnitude Ranking Works}
The cell model explains why ranking by singular value matches the Fisher estimate. Each crossbar column is
quantised with its own scale (\cref{eq:quant}), and the column of component $r$ in the first crossbar holds
$\sigma_r v_r$, so its scale is $s_r=|\sigma_r|\,\|v_r\|_\infty/127$. A fixed cell noise therefore causes a weight error
proportional to $|\sigma_r|$: components with large singular values receive the largest absolute errors.
The Fisher score (\cref{eq:fisher}) multiplies this error size by the loss sensitivity of each stored value.
That it selects essentially the same components as magnitude suggests that the error size varies across
components much more than the sensitivity. The gradient score measures sensitivity to a small change of
$\sigma_r$ alone, ignoring both the error size and the errors on $u_r$ and $v_r$, and so falls behind when
errors are large.

\subsection{Relation to HyFlexPIM}
Our results agree with HyFlexPIM~\cite{song2025hyflexpim} that a hybrid SLC/MLC layout of SVD components
works and that all-MLC storage does not, but not on how to rank components. HyFlexPIM compares its
gradient ranking with magnitude- and rank-based alternatives on MRPC and CoLA at one noise level and reports
differences of about a point or less in favour of the gradient. We find no difference at mild noise
($\sigma=3\%$) and the reverse at high noise, so the findings need not conflict. Three differences in
set-up explain this:
\begin{itemize}
\item \textbf{Noise model.} Under HyFlexPIM's relative weight noise the error of a weight does not depend on
the scale of its crossbar column, so the mechanism above is absent. Its noise level is set from a bit error
rate measured on real MLC cells~\cite{fan2024mlc}, which is itself a per-cell quantity.
\item \textbf{Noise range.} HyFlexPIM evaluates a single noise level. In our sweeps the rankings differ
only at $\sigma\geq5\%$.
\item \textbf{When gradients are taken.} HyFlexPIM records gradients during its 1--3 epochs of
fine-tuning, while ours are taken on the fine-tuned model. We did not repeat their procedure, since magnitude
already matches the Fisher estimate.
\end{itemize}

\subsection{Limitations}
\begin{itemize}
\item \textbf{One error source at a time.} Each non-ideality was varied on its own; their combined effect
on one chip is not measured. The IR-drop model is static, and the ADC range is set by the largest
calibration value rather than a percentile.
\item \textbf{Parameter ranges.} Noise levels, stuck-cell rates and wire resistance were swept from
harmless to failure, not calibrated to a particular device, so we report trends and comparisons between
placements rather than absolute accuracy of a specific chip.
\item \textbf{Known faults.} Fault-aware mapping requires each chip to be tested and mapped once before
programming. Cells that fail later, in use, require testing and remapping again. We did not measure the
time this per-chip mapping takes.
\item \textbf{Cost.} Savings are given in memory cells only; energy, latency and area of the peripheral
circuits, in particular the ADCs, are not evaluated.
\item \textbf{Activation quantisation.} Crossbar inputs are quantised with one scale per tensor. This shifts
SST-2 accuracy by up to about a point between batch sizes, and on MNLI it costs up to 11 points even on an
error-free chip, because the vector between a tier's two crossbars shares one scale across all its
components (\cref{sec:mnli}). All comparisons use identical batches, but finer scales, per token or per
group of components, are needed before absolute accuracies on harder tasks can be trusted.
\item \textbf{Training settings.} Each training variant was run once, with one training chip ($\sigma=5\%$,
0.5\% stuck cells) and 2 epochs; other settings may trade error-free accuracy against robustness
differently.
\item \textbf{Scope.} The full study is on SST-2 with two models; on MNLI only placement and mapping were
run, without training or the other non-idealities. The SVD step itself costs about two points,
and DistilBERT's fine-tuning would benefit from early stopping.
\end{itemize}
